% MIN 6:30 à 7:15. Partir du constat de Sébastien avant d'énoncer la mission.
\begin{frame}{Du copier-coller à des composants réutilisables}
    Des chercheurs reprenaient le code de Sébastien par copier-coller, sans
    toujours le comprendre.

    Mon audit initial et mes reproductions centralisées de zéro m'ont confronté
    à ces difficultés.

    \begin{block}{La mission initiale du stage}
        Réunir les outils nécessaires pour reproduire et comparer ces expériences.
    \end{block}
\end{frame}

% MIN 7:15 à 8:00. Préciser mon périmètre, puis annoncer les trois couches.
\begin{frame}{Krum : une bibliothèque en trois couches}
    Krum est une bibliothèque Python fondée sur \textbf{PyTorch}, un outil
    de calcul et d'entraînement des modèles.

    \begin{tabular}{@{}ll@{}}
        \textbf{Primitives} & Briques de calcul.\\[0.3em]
        \textbf{Simulations} & Exécution des protocoles.\\[0.3em]
        \textbf{Orchestration} & Gestion des campagnes d'expériences.
    \end{tabular}

    \begin{block}{Mes contributions dans ce projet collectif}
        J'ai développé des agrégateurs et travaillé sur le composant \code{Model},
        les architectures de modèles et les simulations centralisés. Le composant \code{Model}
        donne accès aux paramètres sans les recopier.

        J'ai aussi entièrement rédigé la documentation et les tutoriels.
    \end{block}

    \medskip
    Sébastien Rouault et Peva Blanchard ont développé l'orchestrateur.
\end{frame}
